\subsection*{3.1 基本内容}

导数的符号与函数的单调性相联系。若 $f$ 在 $[a,b]$ 上连续、在 $(a,b)$ 内可导，则 $f'>0$ 蕴含严格递增，$f'\geq0$ 蕴含单调递增，$f'\equiv0$ 蕴含 $f$ 为常数。这些结论也可用单侧导数表述。

凸性与导数的单调性相联系。若 $f$ 在开区间上可导，则 $f$ 凸当且仅当 $f'$ 单调递增；若 $f$ 二阶可导，则等价于 $f''\geq0$。严格凸不能简单地等同于 $f''>0$。凸函数在内点具有左右导数，且 $f'_-(x)\leq f'_+(x)$；其图像在每一内点都有支撑线，即存在 $\mu$ 使
\[
f(y)\geq f(x)+\mu(y-x).
\]

\subsection*{3.2 利用导数研究函数的变化}

\ex{3.2.1}若 $f'(0)>0$，能否断定 $f$ 在 $0$ 附近单调递增？若在区间 $(a,b)$ 内处处有 $f'>0$，又能得出什么结论？
\sol 由导数定义只能直接推出点态结论：存在 $\delta>0$，使得 $0<h<\delta$ 时，$f(h)>f(0)>f(-h)$。若再假定 $f'$ 在 $0$ 处连续，则它在某邻域内为正，因而 $f$ 在该邻域内严格递增。对于整个 $(a,b)$，只需 $f'(x)>0$ 处处成立，便可由 Lagrange 中值定理推出 $f$ 严格递增；不必另设 $f'$ 连续。

\ex{3.2.2}讨论导数有界、函数一致连续及凸性的关系。
\sol 若 $|f'|\leq M$，中值定理给出 $|f(x)-f(y)|\leq M|x-y|$，所以 $f$ 是 Lipschitz 连续的，因而一致连续。若 $f'(x)\to+\infty$（$x\to+\infty$），则固定 $h>0$，有
\[
f(x+h)-f(x)=h f'(\xi_x)\longrightarrow+\infty,
\]
故 $f$ 在 $(0,+\infty)$ 上不一致连续。若 $f$ 在该区间上凸且一致连续，则其递增的导数不可能趋于 $+\infty$；同理也不能趋于 $-\infty$。因此 $f'$ 在无穷远处有界。

\ex{3.2.3}设 $f$ 在 $[0,+\infty)$ 上可微。研究一致连续、导数有界与线性增长之间的关系。
\sol 若 $f'$ 有界，则 $f$ 一致连续且至多线性增长。另一方面，一致连续也蕴含至多线性增长：取 $\delta>0$ 使 $|x-y|<\delta$ 时 $|f(x)-f(y)|<1$，将 $[0,x]$ 分成长度小于 $\delta$ 的小段即可。若再假定 $f$ 凸且满足 $|f(x)|\leq A x+B$，则对任意 $h>0$，
\[
f'(x)\leq\frac{f(x+h)-f(x)}{h};
\]
令 $h\to+\infty$ 可得 $f'(x)\leq A$。由于 $f'$ 单调递增，$f'(+\infty)$ 存在且有限。

\ex{3.2.4}研究 $f(+\infty)$ 与 $f'(+\infty)$ 的关系。
\sol 若 $f(+\infty)$ 为有限数，则由中值定理，存在 $x_n\to+\infty$，使
\[
f'(x_n)=\frac{f(2n)-f(n)}{n}\longrightarrow0.
\]
如果 $f$ 还凸，则 $f'$ 单调，从而 $f'(+\infty)=0$。若 $f'(+\infty)$ 存在且不为零，则 $f$ 必趋向相应的无穷大或无穷小。
